\section{Questions}

\subsection{FFNN}
\paragraph{Activation Functions - Pros and Cons}

\paragraph{SGD with Momentum}

\paragraph{Regularization}

\subsection{CNN}
\paragraph{AlexNet, GoogLeNet, ResNet}
\begin{itemize}
    \item \textbf{AlexNet:} Built for the ImageNet competition, started the deep learning revolution in computer vision. 
    Takes the concept of LeNet and makes it deeper. Uses a series of convolutional and max pooling layers, 
    followed by fully connected layers for the final classification. 
    Shown that the feature maps learned by the convolutional layers can be transferred to other tasks.
    \item \textbf{GoogLeNet:}
    \begin{itemize}
        \item Introduced the Inception module, which allows for multi-scale feature extraction by applying multiple convolutional filters of different sizes in parallel.
        \item Uses global average pooling instead of fully connected layers, reducing the number of parameters and improving generalization.
        \item Further reduces the number of parameters by using 1x1 convolutions to reduce the depth of the feature maps before applying larger convolutions.
        \item To avoid the vanishing gradient problem, GoogLeNet uses auxiliary classifiers during training, which provide additional supervision to the intermediate layers.
    \end{itemize}
    \item \textbf{ResNet:} Went much deeper than previous architectures by introducing identity bypass connections (produce an ensemble of shallower networks).
    If one layer will not be needed, it can be skipped. Since the dimensionality needs to be the same, a 1x1 convolution is used to match the channel dimensions.

    Also uses only smaller 3x3 convolutions, which reduces the number of parameters and computational cost.
\end{itemize}

\paragraph{Inception Module - Why 1x1 Convolutions?}
Allow for multy-scale feature extraction by applying multiple convolutional filters of different sizes in parallel.
By doing this, it may also learn redundant features, which makes the model more robust. 

The 1x1 convolutions are used before applying larger convolutions to reduce the depth of the feature maps, 
reducing also the number of parameters and computational cost.

\subsection{Transformer}
\paragraph{Global vs Local attention}
\begin{itemize}
    \item \textbf{Global Attention:} At every step, the decoder compute the attention weights over all the encoder's hidden states.
    Computationally expensive for long sequences.

    Given that we are at step $t$ in the decoder, we first compute the alignment over all the encoder's hidden states $h_1, h_2, ..., h_n$ as:
    \[ e_{t,i} = a(s_{t-1}, h_i) \]
    where $s_{t-1}$ is the decoder's previous hidden state, and $a$ is a scoring function (e.g., dot product, additive, etc.). 
    
    The attention weights are then computed using the softmax (ensure they sum to 1) function:
    \[ \alpha_{t,i} = softmax(e_{t,i}) = \frac{exp(e_{t,i})}{\sum_{j=1}^{n} exp(e_{t,j})} \]
    
    Finally, the context vector $c_t$ is computed as a weighted sum of the encoder's hidden states:
    \[ c_t = \sum_{i=1}^{n} \alpha_{t,i} h_i \]

    \item \textbf{Local Attention:} Instead of attending to all encoder hidden states, the decoder attends to a fixed-size window of encoder hidden states around a predicted position. 
    This reduces computational complexity and allows for more focused attention on relevant parts of the input sequence.

    Given that we are at step $t$ in the decoder, we first predict a position $p_t$ in the encoder's hidden states to attend to. 
    This is typically done using a separate neural network that takes the decoder's previous hidden state $s_{t-1}$ as input and outputs a predicted position $p_t$.
    \[ p_t = T_x \cdot \sigma(v_p^T tanh(W_p s_{t-1})) \]
    where $T_x$ is the length of the input sequence, $v_p$ and $W_p$ are learnable parameters, and $\sigma$ is the sigmoid function.
    This approach allow for backpropagation through the predicted position (the predicted position is differentiable with respect to the model parameters).
    
    To enforce the locality constraint, a Gaussian distribution is applied to the attention weights, centered around $p_t$:
    \[ \alpha_{t,i} = softmax(e_{t,i}) \cdot exp\left(-\frac{(i - p_t)^2}{2\sigma^2}\right) \]
\end{itemize}

\paragraph{Self-Attention}
Given the computed embeddings of the input sequence $x_1, x_2, ..., x_n$, self-attention allows each token to attend to all other tokens in the sequence.
\begin{itemize}
    \item \textbf{Query (Q):} $q_i = W^Q x_i$, where $W^Q$ is a learnable weight matrix.
    The query vector represents what the current token is looking for in the other tokens.
    \item \textbf{Key (K):} $k_i = W^K x_i$, where $W^K$ is a learnable weight matrix.
    The key vector represents the content of the current token that can be matched against the queries of other tokens.
    \item \textbf{Value (V):} $v_i = W^V x_i$, where $W^V$ is a learnable weight matrix.
    The value vector represents the information that will be passed to the next layer if the current token is attended to.
\end{itemize}

Given the queries, keys, and values vectors, the attention scores for input token $i$ attending to all other tokens $j$ are computed as:
\[
e_{ij} = q_i \cdot k_j^T
\]
These scores are then rescaled by the square root of the dimension of the key vectors $d_k$ to prevent large dot product values:
\[
\hat{e}_{ij} = \frac{e_{ij}}{\sqrt{d_k}}
\]
The attention weights are then computed using the softmax function:
\[
\alpha_{ij} = softmax(\hat{e}_{ij}) = \frac{exp(\hat{e}_{ij})}{\sum_{j=1}^{n} exp(\hat{e}_{ij})}
\]
Finally, the output of the self-attention mechanism for token $i$ is computed as a weighted sum of the value vectors:
\[
z_i = \sum_{j=1}^{n} \alpha_{ij} v_j
\]

Similarly, the self-attention mechanism can be defined in matrix form for the entire input sequence as:
\[
\mat{Z} = softmax\left(\frac{\mat{Q} \mat{K}^T}{\sqrt{d_k}}\right) \mat{V}
\]

\subsection{GAN}
\begin{enumerate}
    \item draw the architecture of a GAN.
    \item Write the loss function of a GAN and explain what is the problem with the original loss function? (min-max game, mode collapse, training instability).
    How can we fix it? (Switch to max-max or min-min, LSGAN)
    \item What is Mode Collapse? How can we fix it? (Minibatch discrimination, LSGAN)
    \item Can you create a GAN with some conditional information? 
\end{enumerate}
